\documentclass[a4paper,12pt]{article}

\usepackage[T1]{fontenc}
\usepackage[italian,english]{babel}
\usepackage[utf8]{inputenc}
\usepackage{graphicx}
\usepackage{float}
\usepackage{subcaption}
\usepackage{geometry}
\usepackage{fancyhdr}
\usepackage{xcolor}
\usepackage{mdframed}
\usepackage{enumitem}
\usepackage{hyperref}
\usepackage{booktabs}
\usepackage{array}
\usepackage{multicol}

% ── Brand Colors ─────────────────────────────────────────────────────────────
\definecolor{nutrigreen}{RGB}{46,125,50}
\definecolor{nutrilightgreen}{RGB}{232,245,233}
\definecolor{nutridark}{RGB}{28,28,28}
\definecolor{nutrigray}{RGB}{117,117,117}
\definecolor{nutriaccent}{RGB}{56,142,60}

% ── Page Margins ────────────────────────────────────────────────────────────
\setlength{\headheight}{30pt}
\geometry{left=2.5cm, right=2.5cm, top=2.8cm, bottom=2.5cm}

% ── Header ─────────────────────────────────────────────────────────────
\pagestyle{fancy}
\fancyhf{}
\fancyhead[L]{\small\color{nutrigray}Section 4 --- Application User Guide}
\fancyhead[R]{\small\color{nutrigray}\thepage}
\renewcommand{\headrulewidth}{0.4pt}
\renewcommand{\headrule}{\hbox to\headwidth{\color{nutrigreen}\leaders\hrule height \headrulewidth\hfill}}

% ── Step Counter ───────────────────────────────────────────────────────────
\newcounter{passo}[subsection]
\newcommand{\passaggio}[1]{%
    \stepcounter{passo}%
    \subsubsection*{\color{nutriaccent}Step \thesubsection.\thepasso\ ---\ #1}%
}

% ── Green Info Box ─────────────────────────────────────────────────────
\newmdenv[
  backgroundcolor=nutrilightgreen,
  linecolor=nutrigreen,
  linewidth=1.5pt,
  roundcorner=5pt,
  innertopmargin=8pt,
  innerbottommargin=8pt,
  innerleftmargin=12pt,
  innerrightmargin=12pt,
  skipabove=8pt,
  skipbelow=8pt
]{infobox}

% ── Green Bold Label ─────────────────────────────────────────────────
\newcommand{\labelgreen}[1]{\medskip\noindent{\color{nutriaccent}\textbf{#1}}\par\smallskip}

% ── Italic Tip ──────────────────────────────────────────────────────
\newcommand{\suggerimento}[1]{%
  \smallskip\noindent\textit{\color{nutrigray}#1}\par\smallskip%
}

% ── Variable for Images Folder (useful for translations) ───────────────────
\newcommand{\imgdir}{immagini_ing}

% ── Centered Image with Caption ─────────────────────────────────────────
\newcommand{\schermata}[3][0.38\textwidth]{%
  \begin{figure}[H]
    \centering
    \includegraphics[width=#1]{\imgdir/#2}
    \caption{\small\color{nutrigray}#3}
  \end{figure}%
}

% ── Two Side-by-Side Images ──────────────────────────────────────────────────
\newcommand{\dueschermate}[4]{%
  \begin{figure}[H]
    \centering
    \begin{subfigure}[b]{0.42\textwidth}
      \centering
      \includegraphics[width=\textwidth]{\imgdir/#1}
      \caption{\small\color{nutrigray}#2}
    \end{subfigure}
    \hfill
    \begin{subfigure}[b]{0.42\textwidth}
      \centering
      \includegraphics[width=\textwidth]{\imgdir/#3}
      \caption{\small\color{nutrigray}#4}
    \end{subfigure}
  \end{figure}%
}

% ── Three Side-by-Side Images ──────────────────────────────────────────────────
\newcommand{\treschermate}[6]{%
  \begin{figure}[H]
    \centering
    \begin{subfigure}[b]{0.30\textwidth}
      \centering
      \includegraphics[width=\textwidth]{\imgdir/#1}
      \caption{\small\color{nutrigray}#2}
    \end{subfigure}
    \hfill
    \begin{subfigure}[b]{0.30\textwidth}
      \centering
      \includegraphics[width=\textwidth]{\imgdir/#3}
      \caption{\small\color{nutrigray}#4}
    \end{subfigure}
    \hfill
    \begin{subfigure}[b]{0.30\textwidth}
      \centering
      \includegraphics[width=\textwidth]{\imgdir/#5}
      \caption{\small\color{nutrigray}#6}
    \end{subfigure}
  \end{figure}%
}

% ─────────────────────────────────────────────────────────────────────────────
\begin{document}

% ════════════════════════════════════════════════════════════════════════════
%  SECTION HEADER
% ════════════════════════════════════════════════════════════════════════════
\begin{center}
  {\Huge\bfseries\color{nutrigreen} Section 4}\\[6pt]
  {\Large\bfseries\color{nutridark} Application User Guide}\\[4pt]
  {\normalsize\itshape\color{nutrigray} User Manual --- step-by-step instructions}
\end{center}

\medskip
\noindent\rule{\textwidth}{1.5pt}{\color{nutrigreen}}
\medskip

\noindent This chapter illustrates the main features of \textbf{NutriApp}
through detailed instructions and screenshots taken directly from the application.
Each section is divided into progressive steps to facilitate understanding even
for less experienced users.

\medskip
\begin{infobox}
  \textbf{\color{nutrigreen}Section Index}
  \begin{enumerate}[leftmargin=1.4em, itemsep=2pt]
    \item Login and Registration
    \item Navigation Bar
    \item Adding Food
    \item Setting Goals
    \item Viewing Progress
    \item Account Settings
  \end{enumerate}
\end{infobox}

\newpage

% ════════════════════════════════════════════════════════════════════════════
%  1. LOGIN AND REGISTRATION
% ════════════════════════════════════════════════════════════════════════════
\section{Login and Registration}

Upon opening the application for the first time, the welcome screen is displayed,
from which you can access an existing account or create a new one.
NutriApp supports three login methods: email and password credentials,
authentication via Google account, and authentication via Apple ID.
You can also explore the app in guest mode without creating an account.

% ─── 1.1 ─────────────────────────────────────────────────────────────────────
\subsection{Welcome and Login Screen}

\passaggio{Open the application}

\schermata[0.40\textwidth]{login.jpeg}{Welcome screen of \textbf{NutriApp}:
  at the top is the green logo and the subtitle \emph{Log in to manage your meals,
  macros, and micros all in one app.} The central white panel
  shows the \textbf{Email} and \textbf{Password} fields (with an eye icon to
  show/hide characters), the green \textbf{LOG IN} button, and the
  link \emph{Continue without logging in}. Below the \emph{or} separator are the
  \textbf{Log in with Google} and \textbf{Log in with Apple} buttons, and at the bottom
  is the link \emph{Don't have an account? Sign up}.}

\labelgreen{Description}
When you first open NutriApp, the login screen appears with the app's name
in dark green and a subtitle briefly describing its function.
The central login panel has two fields:

\begin{itemize}[leftmargin=1.4em, itemsep=2pt]
  \item \textbf{Email} --- email address associated with the account,
        indicated by the envelope icon to the left of the field
  \item \textbf{Password} --- account password; touch the crossed-out eye
        icon to the right of the field to toggle between showing the characters
        in plain text or masked
\end{itemize}

After filling out both fields, press the green \textbf{Log In} button
to complete the authentication and enter the application.

\labelgreen{Continue without logging in}
By tapping the link \emph{Continue without logging in}, you can explore the
application in guest mode, without creating or entering credentials.
In this mode, data is not synced to the cloud and is not
accessible from other devices. It is recommended to create an account so you
do not lose entered data.

\suggerimento{If you forget your password, you can recover it via the
recovery procedure available on the official NutriApp website: the
system sends a reset link to the email address associated with the account.}

% ─── 1.2 ─────────────────────────────────────────────────────────────────────
\subsection{Log in with Google}

\passaggio{Select quick login via Google}

\dueschermate{login.jpeg}{Login screen: at the bottom of the white panel is the
  \textbf{Log in with Google} button with the G logo on the left. Tapping it
  opens the selector for Google accounts present on the device.}{login_google.jpeg}{%
  System window \textbf{Choose an account}: shows all Google accounts
  saved on the device with profile picture and email address (obscured
  for privacy). At the bottom is the option \emph{Add another account} to add
  an account not yet present, and the note that Google will share your name,
  email, and profile picture with NutriApp.}

\labelgreen{Description}
Tap \textbf{Log in with Google} on the login screen to start
quick authentication. The operating system automatically opens the Google
account selection window, which shows all Gmail accounts already
configured on the device. Each entry in the list shows the profile picture
and the email address of the account.

\labelgreen{How to proceed}
\begin{itemize}[leftmargin=1.4em, itemsep=2pt]
  \item \textbf{Tap your account} in the list to select it and
        authorize NutriApp to access your name, email, and profile picture.
        The login is completed automatically without further steps.
  \item \textbf{Tap \emph{Add another account}} if the desired account
        is not in the list: the standard Google login flow will open to
        add a new account to the device.
\end{itemize}

\labelgreen{Next step}
After selecting the account, the application authenticates the user and goes
straight to the main screen. If this is the first login with that Google account,
the initial profile setup flow will be started.

\suggerimento{Logging in with Google is the fastest method: it does not require
you to remember a separate password and updates automatically if you change
your Google password. It is recommended for those who already use Google on other services.}

% ─── 1.3 ─────────────────────────────────────────────────────────────────────
\subsection{Create a New Account}

\passaggio{Open the registration form --- Step 1}

\schermata[0.40\textwidth]{registra_step1.jpeg}{Screen \textbf{Registration}
  --- Step 1 of 3: the progress indicator at the top shows the first dash
  as green (active) and two as gray (to be completed). In the center is the placeholder avatar with
  the green camera icon to add an optional photo. The form fields are:
  \textbf{First Name}, \textbf{Last Name}, \textbf{Email}, \textbf{Password}, and
  \textbf{Confirm Password}. At the bottom is the green \textbf{SIGN UP} button.}

\labelgreen{Description}
Tap the green link \emph{Don't have an account? Sign up} on the login screen
to open the three-step registration wizard, displayed
by the progress indicator at the top (green dash = current step,
gray dashes = upcoming steps).

\textbf{Step 1 --- Personal Data and Credentials.}
Fill out the following fields:

\begin{itemize}[leftmargin=1.4em, itemsep=2pt]
  \item \textbf{First Name} and \textbf{Last Name} --- the full name that will be
        displayed in the app and in the greeting on the main screen
  \item \textbf{Email} --- the email address that will act as the
        unique identifier for the account; it must be valid because
        it will be used for verification via an OTP code in the next step
  \item \textbf{Password} --- choose a strong password of at least 8
        characters; it is recommended to include uppercase letters, lowercase letters,
        numbers, and special characters
  \item \textbf{Confirm Password} --- re-enter the same password to
        prevent typos; the two fields must match
        exactly to proceed
\end{itemize}

\textbf{Profile Photo (optional).}
Tap the green camera icon superimposed on the gray avatar to take
a photo or choose one from the gallery. The photo is optional and can be
added or changed at any time later from the
\emph{My Profile} section in the settings.

\labelgreen{Next step}
Check that all required fields are correctly filled out,
then press the green \textbf{SIGN UP} button to proceed to the next
step.

\passaggio{Verify your email address --- Step 2}

\schermata[0.40\textwidth]{registra_step2.jpeg}{Screen \textbf{Registration}
  --- Step 2 of 3: the indicator shows the first two dashes as green and the last
  as gray. In the center is an envelope icon with a green checkmark and the title
  \textbf{Verify your email}. The text indicates that an
  OTP code has been sent to the entered email address (partially obscured).
  The \textbf{OTP Code} field awaits input. The link
  \emph{Send again} allows you to request a new one.
  At the bottom is the green \textbf{CHECK EMAIL} button.}

\labelgreen{Description}
After pressing \textbf{Continue} in the first step, the application automatically
sends an OTP (One-Time Password) code to the entered email address.
The screen shows an envelope icon with a green checkmark and
indicates to which address the code was sent.

\labelgreen{How to complete verification}
\begin{itemize}[leftmargin=1.4em, itemsep=2pt]
  \item Open the inbox of the email address used during
        registration and look for the NutriApp verification email
  \item Note or copy the 6-digit OTP code found in the email
  \item Return to the app and enter the code in the \textbf{OTP Code} field
  \item Press \textbf{CHECK EMAIL} to confirm the email address
        and proceed to the final step
\end{itemize}

\labelgreen{Code not received}
If the code does not arrive within a few minutes, check your provider's
\emph{Spam} or \emph{Junk} folder. Alternatively,
tap the link \emph{Send again} to request a new OTP code.
The link becomes active after a short waiting period to prevent
rapid multiple requests.

\suggerimento{The OTP code is valid for a limited time (usually
10--15 minutes). If you wait too long before entering it, you will
need to request a new one via the \emph{Send again} link.}

\passaggio{Complete registration --- Step 3}

\schermata[0.40\textwidth]{registra_step3.jpeg}{Screen \textbf{Registration}
  --- Step 3 of 3: all three dashes on the indicator are green, signaling
  the completion of the process. In the center is a green circle with a white checkmark and
  the title \textbf{Great job!} with the subtitle \emph{You are ready to
  start your journey with NutriApp.} At the bottom is the green
  \textbf{GO TO HOMEPAGE} button.}

\labelgreen{Description}
After successful verification of the OTP code, the app shows the confirmation screen.
All three dashes of the indicator turn green. The
message \textbf{Great job!} confirms that the account has been created and
verified correctly. Press \textbf{GO TO HOMEPAGE} to enter
the application.

\begin{infobox}
  \textbf{\color{nutrigreen}Summary of Login Methods}

  \smallskip
  \begin{tabular}{@{}lll@{}}
    \toprule
    \textbf{Method}       & \textbf{Requirement}             & \textbf{Ideal for} \\
    \midrule
    Email + Password      & NutriApp Account                 & Those who prefer dedicated credentials \\
    Log in with Google    & Gmail Account on device          & Quick access without password \\
    Log in with Apple     & Apple ID (iOS only)              & iPhone/iPad users \\
    Without logging in    & None                             & Quick exploration of the app \\
    \bottomrule
  \end{tabular}
\end{infobox}

\newpage

% ════════════════════════════════════════════════════════════════════════════
%  2. NAVIGATION BAR
% ════════════════════════════════════════════════════════════════════════════
\section{Navigation Bar}

The navigation bar is the persistent element of the interface: it appears at the
bottom of every screen and allows you to instantly move between the
five main areas of the application with a single tap.

% ─── 2.1 ─────────────────────────────────────────────────────────────────────
\subsection{View the Navigation Bar}

\passaggio{Open the application}

\schermata[0.38\textwidth]{homepage.jpg}{Main screen \textbf{NutriApp}:
  daily caloric summary with circular graph, macronutrient bars
  (Carbs, Protein, Fat) and the four meal time slots. The navigation bar
  at the bottom with the icons NutriApp, Calendar, Add, Recipes,
  and Charts.}

\labelgreen{Description}
When you open the application, the main screen is shown with a
personalized greeting at the top, the current date navigable with the \texttt{<} and \texttt{>} arrows,
and a green circular graph indicating the calories remaining for the day.
To the right of the graph are horizontal colored bars for Carbs (orange),
Protein (green), and Fat (pink/red), each showing the consumed value compared to the goal.
The four time slots --- Breakfast, Lunch, Dinner, and Snack --- are
accessible via the \textbf{+} button next to each one.

\labelgreen{Bar Elements}
\begin{itemize}[leftmargin=1.4em, itemsep=2pt]
  \item \textbf{NutriApp} (house icon) --- main screen with caloric
        summary, macronutrients, and daily meal sections
  \item \textbf{Calendar} (calendar icon) --- monthly view with a green/red color
        code for days with or without logged meals
  \item \textbf{Add} (central green \textbf{+} button) --- quick access
        to the food entry screen
  \item \textbf{Recipes} (book icon) --- list of saved recipes
  \item \textbf{Charts} (line chart icon) --- nutritional statistics
        with daily, weekly, monthly, and yearly views
\end{itemize}

\suggerimento{The active icon is highlighted in green on the bottom bar.
The green \textbf{+} button in the center is always available to quickly add
a food item: it's the fastest route for daily logging.}

\newpage

% ════════════════════════════════════════════════════════════════════════════
%  3. ADDING FOOD
% ════════════════════════════════════════════════════════════════════════════
\section{Adding Food}

The application offers multiple ways to log meals: searching the global database,
selecting from saved foods, adding an already created recipe, and
scanning the product's barcode.

% ─── 3.1 ─────────────────────────────────────────────────────────────────────
\subsection{Open the Entry Screen}

\passaggio{Go to the main screen}

\schermata[0.38\textwidth]{homepage.jpg}{HomePage: circular graph with remaining
  calories, macronutrient bars, and the four time slots each with
  the \textbf{+} button.}

\labelgreen{Next step}
Press the \textbf{+} next to the desired time slot (Breakfast, Lunch,
Dinner, or Snack) or press the large green \textbf{+} button in the center of the
navigation bar to access the \emph{Add} screen.

% ─── 3.2 ─────────────────────────────────────────────────────────────────────
\subsection{Choose the Entry Method}

\passaggio{Select the food source}

\dueschermate{pasto_dropdown.jpg}{Screen \textbf{Manual}: dropdown menu
  \emph{Select Meal} and fields for Food Name, Weight (g), and Macronutrients
  (Calories, Carbs, Protein, Fat, Fiber, Sugars).}{pasto_recenti.jpg}{%
  Screen \textbf{Add} with active tab \emph{Add Food}: search bar
  with barcode icon. Below is the link \emph{Can't find the food?
  Register it yourself!}}

\labelgreen{Description}
The \emph{Add} screen has three tabs at the top:

\begin{itemize}[leftmargin=1.4em, itemsep=2pt]
  \item \textbf{Add Food} --- search the global database or scan
        the barcode
  \item \textbf{Saved Recipes} --- recipes created by the user with the total
        calories of each
  \item \textbf{Saved Foods} --- previously manually registered
        foods
\end{itemize}

\passaggio{Navigate between Saved Recipes and Saved Foods tabs}

\dueschermate{pasto_ricette.jpg}{Tab \textbf{Saved Recipes}: each recipe
  shows its name and total calories. At the bottom is the link
  \emph{Insert new recipe}.}{pasto_preferiti.jpg}{Tab \textbf{Saved
  Foods}: search bar; if there are no saved foods, the message
  \emph{No saved food} appears.}

\suggerimento{Manually entered foods are automatically saved
in the \emph{Saved Foods} section and will be available for future entries
without having to re-enter them from scratch.}

% ─── 3.3 ─────────────────────────────────────────────────────────────────────
\subsection{Manual Food Entry}

\passaggio{Fill out the entry form}

\dueschermate{registra_alimento.jpg}{Form \textbf{Ingredient}: fields
  \emph{Food Name} and \emph{Weight (g)}, followed by the
  \textbf{Macronutrients} section with Carbs, Protein, Fat, Fiber,
  Sugars, and Water.}{nutrienti_micro1.jpg}{Section \textbf{Ingredient} with
  expandable sections \emph{Fat Details}, \emph{Vitamins}, and
  \emph{Minerals} for complete nutritional tracking.}

\labelgreen{Description}
Tap \emph{Register it yourself!} to open the manual entry form.
Fill out at least the mandatory fields in the \textbf{Macronutrients} section by
reading the values from the product packaging (\emph{Nutritional values
per 100\,g} label). For complete tracking, you can expand
the \textbf{Fat Details}, \textbf{Vitamins}, and \textbf{Minerals} sections.


% ─── 3.4 ─────────────────────────────────────────────────────────────────────
\subsection{Entry via Barcode Scan}

\passaggio{Start the scanning camera}

\schermata[0.38\textwidth]{scansione_barcode.jpg}{Scanning screen
  \textbf{barcode}: the camera focuses on the back of a package.
  The green frame delimits the recognition area. At the top are the controls
  \textbf{FLASH ON} and \textbf{CANCEL}.}

\labelgreen{Description}
Tap the barcode icon to the right of the search bar in the
\emph{Add Food} tab. Frame the barcode within the
green box: recognition happens automatically as soon as the
code comes into focus, without pressing any buttons.

\suggerimento{In poorly lit environments, activate the flash by tapping
\textbf{FLASH ON} at the top right. Hold the phone about 15--20\,cm
from the package for optimal results.}

% ─── 3.5 ─────────────────────────────────────────────────────────────────────
\subsection{Create a New Recipe}

\passaggio{Open the form and add ingredients}

\treschermate{nuova_ricetta.jpg}{Form \textbf{Create Recipe}: fields
  \emph{Recipe Name}, \emph{Portions}, and \emph{Procedure / Notes}.
  The Ingredients section shows \emph{0 added}. At the bottom are the buttons
  \textbf{+ Ingredients} and \textbf{Save}.}{nuova_ricetta_recenti.jpg}{%
  Screen \textbf{Add Ingredient} with tab \emph{Add Food}:
  search bar with barcode. Below is the link
  \emph{Can't find the food? Register it yourself!}}{nuova_ricetta_preferiti.jpg}{%
  Tab \emph{Saved Foods}: list of already registered foods or
  the message \emph{No saved food} if empty.}

\labelgreen{Description}
From the \emph{Recipes} section, tap the green \textbf{+} button to open
the \emph{Create Recipe} form. Fill out the name, number of portions, and optional notes,
then press \textbf{+ Ingredients} to add each
component. Each added ingredient appears in the list with edit
and delete icons. When finished, press \textbf{Save}.

% ─── 3.6 ─────────────────────────────────────────────────────────────────────
\subsection{View Meal Details}

\passaggio{Open the time slot detail screen}

\schermata[0.38\textwidth]{dettaglia_pasto.jpg}{Screen \textbf{Breakfast}: orange banner
  with icon and food counter, date navigation, \textbf{+} button
  to add foods. The \emph{Nutritional Summary} shows progress bars
  for Calories, Carbs, Protein, and Fat. At the bottom is the list of
  \emph{Consumed Foods} with edit and delete icons.}

\labelgreen{Description}
Tap a time slot on the main screen to open the detailed summary
with the colored banner, the nutritional summary with progress bars,
and the list of consumed foods with the edit (blue pencil)
and delete (red trash can) icons.

\newpage

% ════════════════════════════════════════════════════════════════════════════
%  4. SETTING GOALS
% ════════════════════════════════════════════════════════════════════════════
\section{Setting Goals}

Goals define the reference values for calories, macronutrients,
and weight that the app uses to calculate progress. Setting realistic goals
is essential for obtaining meaningful tracking data.

% ─── 4.1 ─────────────────────────────────────────────────────────────────────
\subsection{Access and Configure Goals}

\passaggio{Open Settings and select My Goals}

\dueschermate{impostazioni.jpg}{Screen \textbf{Settings}: sections for Account
  (My profile, My goals, Notifications, Language and
  country), Support (Help, Privacy policy, About us), and Actions (Report an
  issue, Log out).}{obiettivi.jpg}{%
  Screen \textbf{My goals}: \emph{Weight} section with \emph{Current
  weight} (60.0\,kg) and \emph{Goal weight} (70.0\,kg). Section
  \emph{Macronutrients} with \emph{Calorie goal} (2000.0\,kcal),
  \emph{Carbs} (130.0\,g), \emph{Protein} (45.0\,g), and
  \emph{Fat} (36.0\,g). At the bottom are the expandable sections
  \emph{Macro Details}, \emph{Vitamins}, and \emph{Minerals}.}

\labelgreen{Description}
Tap the gear icon in the top right to open settings,
then select \textbf{My goals}. The screen gathers the parameters
into two sections:

\begin{itemize}[leftmargin=1.4em, itemsep=2pt]
  \item \textbf{Weight} --- the \emph{Current weight} is the starting point for
        calculating needs; the \emph{Goal weight} is the target
        to reach (can be higher or lower than the current weight)
  \item \textbf{Macronutrients} --- the \emph{Calorie goal} defines the daily
        caloric limit; the Carbs, Protein, and Fat fields set
        the limits in grams
\end{itemize}

\labelgreen{Next step}
Enter the desired values and press \textbf{Save}. The changes are
applied immediately throughout the application. The green confirmation banner
appears at the bottom of the screen.

\begin{infobox}
  \textbf{\color{nutrigreen}How to choose the right values:}
  a common starting point is carbs 45--65\%, protein 10--35\%,
  fat 20--35\% of total daily calories. Consult a
  nutrition professional for personalized values.
\end{infobox}

\newpage

% ════════════════════════════════════════════════════════════════════════════
%  5. VIEWING PROGRESS
% ════════════════════════════════════════════════════════════════════════════
\section{Viewing Progress}

The application offers four time views to analyze the caloric and
nutritional trends: daily, weekly, monthly, and yearly,
accessible from the \textbf{Charts} section in the navigation bar.

% ─── 5.1 ─────────────────────────────────────────────────────────────────────
\subsection{Access Progress Data}

\passaggio{Open the Charts section}

\schermata[0.38\textwidth]{storico_giornaliero.jpg}{Screen \textbf{Charts}
  with active \emph{Daily} tab: green box \emph{Nutritional
  statistics} with logged items, tracked days, and date of last
  entry. Below are the filter selector, the four time tabs, and the red box
  with the total calories.}

\labelgreen{Description}
Tap the \textbf{Charts} icon in the navigation bar. At the top appears
the green box with the global summary, then the \textbf{Filters} panel
with the selected nutrient (default: Calories) and the date range.
The four tabs allow you to change the granularity of the view.

% ─── 5.2 ─────────────────────────────────────────────────────────────────────
\subsection{Analyze the Charts}

\passaggio{Bar Chart --- Calories Over Time}

\schermata[0.40\textwidth]{grafici_barre.jpeg}{Screen \textbf{Charts} with
  active \emph{Daily} tab. At the top is the \textbf{Export PDF report} panel
  with the \textbf{Print PDF} (green) and \textbf{Share PDF} buttons. The
  \textbf{Calories over time} chart shows three red vertical bars for
  April 28, 29, and 30: 779\,kcal, 143\,kcal, and 377\,kcal. The scale
  on the vertical axis goes from 0.0 to 779\,kcal. Below starts the
  \emph{Macronutrient distribution} section.}

\labelgreen{Description}
The \textbf{Calories over time} bar chart shows the daily
caloric intake for the selected period. Each red vertical bar represents
one day, with the value in kcal indicated above. The height is proportional
to the total calories: days with few logged foods produce shorter
bars (like April 29 with 143\,kcal), days with more meals produce
taller bars (like April 28 with 779\,kcal). The vertical axis automatically adapts
to the maximum value of the period.

\labelgreen{Export PDF Report}
Above the chart, the \textbf{Export PDF report} panel offers two buttons:
\begin{itemize}[leftmargin=1.4em, itemsep=2pt]
  \item \textbf{Print PDF} (green) --- generates the document and opens the operating
        system's print window
  \item \textbf{Share PDF} (white with border) --- generates the document and
        opens the share panel to send it via email, WhatsApp, or
        any other installed app
\end{itemize}

\passaggio{Pie Chart --- Macronutrient Distribution}

\schermata[0.40\textwidth]{grafici_torta.jpeg}{Screen \textbf{Charts} ---
  \textbf{Macronutrient distribution} section: ring chart with the
  percentage breakdown of macronutrients. The legend indicates Carbs
  44.7\%, Protein 12.0\%, and Fat 43.3\%. The horizontal bars below
  show the absolute values: Carbs 155.6\,g (orange bar),
  Protein 41.9\,g (blue bar), and Fat 67.1\,g (purple bar).}

\labelgreen{Description}
Scrolling down in the \emph{Charts} screen reveals the
\textbf{Macronutrient distribution} section, with two complementary representations:

\begin{itemize}[leftmargin=1.4em, itemsep=2pt]
  \item \textbf{Ring Chart} --- visually shows the proportion among
        the three macronutrients: the golden/orange segment represents
        Carbs, the dark blue one Protein, and the purple one Fat.
        The percentages in the legend on the right allow you to read the exact
        values at a glance
  \item \textbf{Horizontal Bars} --- show both the percentage and the
        absolute value in grams of each macronutrient, making it easy to
        compare with the goals set in the \emph{My goals} section
\end{itemize}

\labelgreen{How to interpret the data}
In the example shown, Fat makes up 43.3\% of the total energy
intake (67.1\,g), a value higher than the recommended threshold
of 20--35\%. Carbs stand at 44.7\% (155.6\,g) and Protein
at 12.0\% (41.9\,g). This information allows you to identify imbalances
in the macronutrient distribution and adjust food choices
in the following days.

\labelgreen{Next step}
Use the \textbf{Filters} panel at the top of the screen to change the
displayed nutrient (e.g., from Calories to Carbs) and the date range.
Tap the \textbf{Weekly}, \textbf{Monthly}, or \textbf{Yearly} tabs to
view the trend over longer periods.

\suggerimento{Comparing the ring chart and the goals in
\emph{My goals} is the most effective way to evaluate if your
diet is aligned with your targets. If one slice is much larger than the others,
it's advisable to rebalance your meals in the following days.}

% ─── 5.3 ─────────────────────────────────────────────────────────────────────
\subsection{Using the Calendar}

\passaggio{Open the Calendar}

\schermata[0.38\textwidth]{calendario.jpg}{Screen \textbf{Calendar} of
  April 2026: days in red without logged meals, the 28th in green with
  a blue border (current day with logged meal), future days in gray.
  At the bottom is the legend: green = Logged meals, red = No meals.}

\labelgreen{Description}
Tap the \textbf{Calendar} icon in the navigation bar. The color
code immediately identifies logged days: green (meal present),
red (no meal), light gray (future days), blue border (current day).
Tap a day to view the detail of the meals for that date.

\newpage

% ════════════════════════════════════════════════════════════════════════════
%  6. ACCOUNT SETTINGS
% ════════════════════════════════════════════════════════════════════════════
\section{Account Settings}

The \emph{Settings} section gathers all personalization options:
personal profile, notifications, units of measure, language, and account management.

% ─── 6.1 ─────────────────────────────────────────────────────────────────────
\subsection{Open Settings}

\passaggio{Access the Settings screen}

\schermata[0.38\textwidth]{impostazioni.jpg}{Screen \textbf{Settings}:
  the three sections Account, Support, and Actions with all available
  options, each with its own icon.}

\labelgreen{Description}
Tap the gear icon in the top right to open settings.

\medskip
\begin{infobox}
  \begin{tabular}{@{}ll@{}}
    \toprule
    \textbf{Option}            & \textbf{Function} \\
    \midrule
    My profile                 & Change first name, last name, height, birth date, gender, and photo \\
    My goals                   & Set target calories, weight, macros, and micronutrients \\
    Notifications              & Meal and hydration reminders with custom times \\
    Language and country       & Language (IT, EN, ZH, AR) and reference country \\
    \midrule
    Help                       & Help center with frequently asked questions \\
    Privacy policy             & Document on data management policies \\
    About us                   & Information about the application and the team \\
    \midrule
    Report an issue            & Direct report to the support team \\
    Log out                    & Log out from the current account \\
    \bottomrule
  \end{tabular}
\end{infobox}

% ─── 6.2 ─────────────────────────────────────────────────────────────────────
\subsection{Modify Account Data}

\passaggio{Update your personal profile}

\schermata[0.38\textwidth]{profilo.jpg}{Screen \textbf{My profile}:
  profile photo with green camera icon, fields \emph{First Name} (Christian),
  \emph{Last Name} (Donzelli), \emph{Height} (175.0), \emph{Date of Birth}
  (11/30/-1) and dropdown menu \emph{Gender} (Male). At the bottom is the
  green \textbf{Save} button.}

\labelgreen{Description}
Tap \textbf{My profile} to modify your personal data:
\begin{itemize}[leftmargin=1.4em, itemsep=2pt]
  \item \textbf{Profile photo} --- tap the green camera icon to take
        a photo or choose one from your gallery
  \item \textbf{First and Last Name} --- the user's personal information
  \item \textbf{Height} --- in centimeters, used for needs calculation
  \item \textbf{Date of Birth} --- used to calculate the basal metabolic rate
  \item \textbf{Gender} --- affects the caloric needs formula
\end{itemize}

Press \textbf{Save} to confirm. The app automatically updates the nutritional
calculations based on the new data.

% ─── 6.3 ─────────────────────────────────────────────────────────────────────
\subsection{Configure Notifications}

\passaggio{Enable reminders}

\dueschermate{notifiche.jpg}{Screen \textbf{Notifications} with active toggles:
  \emph{Meal reminders} with times Breakfast 7:30\,AM, Lunch 12:00\,PM,
  Dinner 7:00\,PM, Snacks 3:00\,PM; \emph{Water reminder} with frequency
  \emph{Every 30 minutes}.}{notifiche_orari.jpg}{Dropdown \emph{Reminder
  frequency} open: \emph{Every 30 minutes} (selected), \emph{Every 1 hour},
  \emph{Every 2 hours}, and \emph{I DECIDE}. At the bottom is the \emph{Alert History}
  with recent notifications.}

\labelgreen{Description}
In the \emph{Notifications} screen, two toggles are available:
\begin{itemize}[leftmargin=1.4em, itemsep=2pt]
  \item \textbf{Meal reminders} --- turning this on displays the four time
        slots; tap a green time to modify it with the system selector
  \item \textbf{Water reminder} --- turning this on displays the frequency menu
        (Every 30 min, Every 1 hour, Every 2 hours, I DECIDE)
\end{itemize}
At the bottom is the \textbf{Alert History} with recent notifications; tap
\textbf{Clear} to delete it.

% ─── 6.4 ─────────────────────────────────────────────────────────────────────
\subsection{Change Language and Country}

\passaggio{Configure localization}

\schermata[0.38\textwidth]{lingua_paese.jpg}{Screen \textbf{Language and Country}: \emph{Language} menu open with \emph{Italian} selected, \emph{English}, \emph{Simplified Chinese}, and \emph{Arabic}. Below is the \emph{Country} section with \emph{Italy}, \emph{England}, \emph{China}, and \emph{Egypt}.}

\labelgreen{Description}
The available languages are Italian, English, Simplified Chinese, and Arabic. The change applies to the interface when the app is restarted.

% ─── 6.5 ─────────────────────────────────────────────────────────────────────
\subsection{Report an Issue}

\passaggio{Send a report to support}

\schermata[0.38\textwidth]{segnala_problema.jpg}{Screen \textbf{Report an
  issue}: text field with placeholder \emph{Describe in detail
  the issue encountered\ldots} and green \textbf{SEND REPORT} button.}

\labelgreen{Description}
Tap \textbf{Report an issue} in the \emph{Actions} section. Fill out the
text field with a detailed description and press \textbf{SEND REPORT}.

\suggerimento{For a faster resolution, include in your report the
device model (e.g., iPhone 14, Samsung Galaxy S23), the version of the
operating system, and the precise description of the steps that led to the problem.}

\bigskip
\noindent\rule{\textwidth}{1pt}{\color{nutrigreen}}
\begin{center}
  \small\itshape\color{nutrigray}--- End of Section 4 ---
\end{center}

\end{document}
